\subsection{指数映射的函子性}

命题 8. 设 G 和 H 是李群芽，h 是从 G 到 H 的态射，\(\phi_{\mathrm{G}}\) 和 \(\phi_{\mathrm{H}}\) 分别是 G 和 H 的指数映射。存在 L(G) 中 0 的一个邻域 V，使得 \(h \circ \phi_{\mathrm{G}}\) 与 \(\phi_{\mathrm{H}} \circ \mathrm{L}(h)\) 在 V 上重合。

给 \(\mathrm{L}(\mathrm{G})\) 和 \(\mathrm{L}(\mathrm{H})\) 赋予分别定义其拓扑且满足 \(\| [x,y]\| \leqslant \| x\| \| y\|\) 对所有 \(x\) 和 \(y\) 成立的范数。可以假设 \(\mathrm{G}\)（分别为 H）是由 \(\mathrm{L}(\mathrm{G})\)（分别为 \(\mathrm{L}(\mathrm{H})\)）定义的李群芽，从而 \(\phi_{\mathrm{G}}\)（分别为 \(\phi_{\mathrm{H}}\)）在 0 的某个邻域上与 \(\mathrm{Id}_{\mathrm{G}}\)（分别为 \(\mathrm{Id}_{\mathrm{H}}\)）重合。另一方面，\(\mathrm{L}(\mathrm{G})\) 中存在 0 的一个对称开邻域 W，使得 \(\mathrm{L}(h)\) 是从李群芽 W 到 H 的态射。根据定理1，\(\mathrm{L}(h)\) 在 0 的某个邻域上与 \(h\) 重合，命题得证。

粗略地说，若借助指数映射在单位元的某个邻域上将 G 和 H 分别与 L(G) 和 L(H) 识别，则从 G 到 H 的每个态射在 0 的某个邻域上都是线性的。

推论 1. 设 G 是李群芽，G' 是 G 的李子群芽，\(\phi\) 是 G 的指数映射。

\begin{enumerate}
\def\labelenumi{(\roman{enumi})}
\item
  存在一个开邻域 \(\mathbf{V}\)，它位于 \(\mathrm{L}(\mathbf{G}')\) 中的 0 附近，使得 \(\phi | \mathbf{V}\) 是从流形 \(\mathbf{V}\) 到 \(e\) 在 \(\mathbf{G}'\) 中的一个开邻域的同构。
\item
  设 \(x \in \mathrm{L}(\mathrm{G})\)。以下条件等价：(a) \(x \in \mathrm{L}(\mathrm{G}')\)；(b) \(\phi(\lambda x) \in \mathrm{G}'\)，当 \(|\lambda|\) 充分小时。
\item
  将命题8应用于从 \(\mathbf{G}'\) 到 \(\mathbf{G}\) 的典范嵌入即可得到，且 (ii) 由 (i) 推出。
\end{enumerate}

推论 2. 设 G 是李群，\(\rho\) 是 G 的解析线性表示，\(\phi\) 是 G 的指数映射。存在 L(G) 中 0 的一个邻域 V，使得

对所有 \(x \in V\) 均成立。

\[
\rho (\phi (x)) = \exp (\mathrm{L} (\rho) x)
\]

这由命题8及第3号的例2推出。

推论 3. 设 G 是李群，\(\phi\) 是 G 的指数映射。

\begin{enumerate}
\def\labelenumi{(\roman{enumi})}
\tightlist
\item
  存在 L(G) 中 0 的一个邻域 V，使得
\end{enumerate}

\[
\operatorname{Ad} (\phi (x)) = \exp \operatorname{ad} x.
\]

对所有 \(x \in V\) 均成立。

\begin{enumerate}
\def\labelenumi{(\roman{enumi})}
\setcounter{enumi}{1}
\tightlist
\item
  若 \(g \in \mathbf{G}\)，则存在 0 的一个邻域 \(\mathbf{W}\)，它位于 \(\mathrm{L}(\mathbf{G})\) 中，使得
\end{enumerate}

\[
g \phi (x) g ^ {- 1} = (\mathrm{Ad} g. x)
\]

对所有 \(x \in W\) 均成立。

\begin{enumerate}
\def\labelenumi{(\roman{enumi})}
\tightlist
\item
  这由推论2及 § 3 第12号命题44推出。
\end{enumerate}

\[
g.
\]

